\documentclass{article}

\usepackage[margin=1in]{geometry}
\usepackage{xcolor}
\usepackage{stepdiff}

\title{stepdiff v1.1.0 Demo}
\author{}
\date{}

\stepdiffsetup{
  display=notes,
  theme=soft,
  layout=lecture,
  highlight=background,
  reason-style=muted
}

\begin{document}

\maketitle

\section*{Algebra expansion}

\[
\begin{stepdiff}[frame=true]
\step{(x+2)(x+3)}
\step[reason={distribute}]{x(x+3)+2(x+3)}
\step[reason={collect terms}]{x^2+5x+6}
\step[reason={final form}, tag=final]{x^2+5x+6}
\end{stepdiff}
\]

\section*{Relation-aware derivation}

\[
\begin{stepdiff}[layout=wide]
\step{a_n \le b_n}
\step[reason={take limits}]{\lim a_n \le \lim b_n}
\step[reason={identify limits}, tag=final]{L \le M}
\end{stepdiff}
\]

\section*{Derivative calculation}

\[
\begin{stepdiff}[display=focus, theme=focus, reason-style=badge]
\step{f(x) = x^3 + 2x}
\step[reason={differentiate}]{f'(x) = 3x^2 + 2}
\step[reason={evaluate at x=2}]{f'(2) = 3\cdot 2^2 + 2}
\step[reason={final value}, tag=final]{f'(2) = 14}
\end{stepdiff}
\]

\section*{Manual control}

\[
\begin{stepdiff}[theme=minimal, highlight=underline]
\step{x^2 - 1 = (x-1)(x+1)}
\step[reason={automatic diff}, diff=auto]{x^2 - 1 = x^2 - 1}
\step[reason={suppress noisy highlight}, diff=false]{x^2 - 1 = x^2 - 1}
\step[reason={emphasize identity}, diff=all, tag=final]{x^2 - 1 = (x-1)(x+1)}
\end{stepdiff}
\]

\section*{Customization hook}

\begingroup
\renewcommand{\SDchanged}[1]{\begingroup\color{red!65!black}#1\endgroup}
\renewcommand{\SDreason}[1]{\normalfont\scriptsize\sffamily\color{black!55}#1}
\[
\begin{stepdiff}[frame=true]
\step{p(x) = x^2 + 2x + 1}
\step[reason={custom highlight style}, tag=final]{p(x) = (x+1)^2}
\end{stepdiff}
\]
\endgroup

\end{document}
